\scanpage{10}
% archive foliation: 9
% Typescript upside down in the scan; reading orientation restored here.
\begin{center}-- 36 --\end{center}
\hfill n$^{\circ}$329

si $U$ est ouvert dans $A$ et saturé pour $R_A$, il résulte de l'hypothèse que $u^{-1}(U)$ est identique à l'ensemble obtenu en saturant $U$ pour $R$ ; comme $u$ est continue, $u^{-1}(U)$ est ouvert dans $X$ (§ 2, n$^{\circ}$1, th. 1) d'où le corollaire, en vertu de la prop. 10.

Exemples. -- 1) Soit $X$ un espace somme d'une famille $(X_i)_{i\in I}$ d'espaces topologiques, identifiés canoniquement à des sous-espaces de $X$, et considérons l'espace quotient $X/R$ obtenu par \uline{recollement} des $X_i$ le long de sous-ensembles $A_{i\unsure{x|\varkappa}}$ au moyen de bijection $h_{i\unsure{x|\varkappa}}$ (n$^{\circ}$4, \uline{Exemple} 1). Soit $f$ l'application canonique $X\to X/R$ ; sa restriction à chaque $X_i$ est par hypothèse une bijection continue de $X_i$ sur le sous-espace $f(X_i)$ de $X/R$ ; cherchons si cette application peut être un \uline{homéomorphisme} quel que soit $i\in I$. Puisque $f(A_{i\unsure{x|\varkappa}})=f(A_{\unsure{x|\varkappa}i})$, une condition \uline{nécessaire} est que $h_{i\unsure{x|\varkappa}}$ soit un homéomorphisme quels que soient $i$ et $\unsure{x|\varkappa}$. Mais cette condition \uline{n'est pas suffisante} en général (exerc.\quad). Elle l'est toutefois dans le cas important où les $A_{i\unsure{x|\varkappa}}$ sont tous \uline{ouverts} ou tous \uline{fermés.} En effet, supposons par exemple que les $A_{i\unsure{x|\varkappa}}$ soient tous fermés ; si $F$ est un ensemble fermé quelconque dans $X_i$, $h_{i\unsure{x|\varkappa}}(F\cap A_{i\unsure{x|\varkappa}})$ est fermé dans $X_{\unsure{x|\varkappa}}$ ; or, pour saturer $F$ pour $R$, il faut prendre la réunion de tous les ensembles $h_{i\unsure{x|\varkappa}}(F\cap A_{i\unsure{x|\varkappa}})$ (pour $\unsure{x|\varkappa}\in I$) ; on obtient ainsi un ensemble fermé dans $X$ en vertu du § 1, n$^{\circ}$5, prop. 4, et notre assertion résulte donc de la prop. 10. Le raisonnement est analogue si les $A_{i\unsure{x|\varkappa}}$ sont ouverts.

$*$ 2) Désignons par $R$ la relation d'équivalence $x\equiv y$ (mod.1) dans la droite numérique \uline{$R$} (n$^{\circ}$4, Exemple 2), par $A$ l'intervalle fermé $[0,1]$ ; $A$ contient un point au moins de toute classe d'équivalence suivant $R$. L'application canonique de $A/R_A$ sur le tore \uline{$T$} est un homéomorphisme ; en effet, soit $F$ un ensemble fermé dans $A$ (donc dans \uline{$R$}) ; pour saturer $F$ pour la relation $R$, il faut prendre la réunion des ensembles fermés $F+n$ (pour $n\in\text{\uline{$Z$}}$), qui forment évidemment une famille localement finie ; leur réunion est donc fermée (§ 1, n$^{\circ}$5, prop. 4), d'où notre
